%-------------------------------------------------------------------------------
% Jiaming Chu — One-page Resume (AI / Research Internship)
%
% Layout adapted from Jake Gutierrez's "Jake's Resume" (MIT License,
% github.com/jakegut/resume), the de-facto standard for ML/SWE internship
% resumes, with the section order of an academic CV (Education -> Research ->
% Publications/Projects -> Experience -> Skills) as used by many AI researchers.
%
% Build:  latexmk -pdf resume.tex      (pdflatex; text stays ATS-parseable)
% Lines marked "% TODO" are placeholders to fill in or delete.
%-------------------------------------------------------------------------------
\documentclass[a4paper,10pt]{article}

\usepackage[T1]{fontenc}
\usepackage[utf8]{inputenc}
\usepackage[sfdefault]{FiraSans}       % clean sans; swap for \usepackage{libertine} for a serif look
\usepackage[margin=0.6in,top=0.5in,bottom=0.5in]{geometry}
\usepackage{titlesec}
\usepackage{enumitem}
\usepackage{tabularx}
\usepackage{xcolor}
\usepackage{fontawesome5}
\usepackage[hidelinks]{hyperref}
\usepackage{microtype}

\definecolor{accent}{HTML}{1F4E79}
\definecolor{subtle}{HTML}{555555}
\hypersetup{colorlinks=true,urlcolor=accent,linkcolor=accent}
\pdfgentounicode=1                      % make the PDF text machine-readable for ATS

\pagestyle{empty}
\setlength{\parindent}{0pt}
\setlength{\tabcolsep}{0pt}
\raggedright

%---------------- Section headings ----------------
\titleformat{\section}
  {\large\bfseries\color{accent}\scshape}{}{0em}{}[\vspace{-6pt}\color{accent}\rule{\linewidth}{0.6pt}]
\titlespacing*{\section}{0pt}{8pt}{5pt}

%---------------- Entry macros ----------------
% \entry{Title}{Right-aligned date}{Subtitle}{Right-aligned location}
\newcommand{\entry}[4]{%
  \begin{tabularx}{\linewidth}{@{}X r@{}}
    \textbf{#1} & #2 \\
    \textit{\small #3} & \textit{\small\color{subtle} #4} \\
  \end{tabularx}\vspace{-3pt}}

% \project{Title}{links}{date}
\newcommand{\project}[3]{%
  \begin{tabularx}{\linewidth}{@{}X r@{}}
    \textbf{#1}\ \ {\small #2} & #3 \\
  \end{tabularx}\vspace{-3pt}}

% \pub{Authors}{Title}{Venue}{links}
\newcommand{\pub}[4]{\item #1. \textbf{#2}. \textit{#3}. {\small #4}}

\newcommand{\me}{\underline{Jiaming Chu}}
\newcommand{\sep}{\ \textbar\ }

\setlist[itemize]{leftmargin=1.2em,itemsep=1pt,topsep=2pt,parsep=0pt,label=\small$\bullet$}

%===============================================================================
\begin{document}

%---------------- Header ----------------
\begin{center}
  {\Huge\bfseries Jiaming Chu}\\[4pt]
  \small
  \faEnvelope\ \href{mailto:jiaming.chu@sjtu.edu.cn}{jiaming.chu@sjtu.edu.cn}\quad
  \faGlobe\ \href{https://chu-jiaming.github.io}{chu-jiaming.github.io}\quad
  \faGithub\ \href{https://github.com/chu-jiaming}{github.com/chu-jiaming}\quad
  \faPhone\ +86 000-0000-0000 % TODO: phone, or delete
\end{center}
\vspace{-6pt}

%---------------- Education ----------------
\section{Education}
\entry{Shanghai Jiao Tong University}{Sep 2025 -- Present}
      {M.Eng. in Electronic Information}{Shanghai, China}
\begin{itemize}
  \item Advisor: TODO \sep GPA: TODO % TODO: fill in or delete this line
\end{itemize}

\vspace{3pt}
\entry{Harbin Institute of Technology}{Sep 2021 -- Jun 2025} % TODO: check start date
      {B.Eng. in IntelliSense Engineering}{Harbin, China}
\begin{itemize}
  \item GPA: TODO \sep Rank: TODO % TODO: fill in or delete this line
\end{itemize}

%---------------- Research Interests ----------------
\section{Research Interests}
Generative models, 3D reconstruction (NeRF / 3D Gaussian Splatting), and world models;
robustness of radiance fields to imperfect inputs such as noisy camera poses, and
evaluation protocols whose metrics measure what we think they measure.

%---------------- Publications ----------------
% Uncomment once you have papers / preprints. \me underlines your own name.
% \section{Publications}
% \begin{itemize}
%   \pub{\me, A. Coauthor, B. Advisor}{Paper Title}{Conference 2027}
%       {[\href{https://arxiv.org/abs/xxxx}{arXiv}] [\href{https://github.com/}{code}]}
% \end{itemize}

%---------------- Research Projects ----------------
\section{Research Projects}
\project{PoseGuard-NeRF: When Does Joint Camera-Pose Optimization Help NeRF?}
  {[\href{https://chu-jiaming.github.io/poseguard-nerf/}{project}]
   [\href{https://github.com/chu-jiaming/PoseGuard-NeRF}{code}]}
  {2026}
\begin{itemize}
  \item Designed a controlled study of \textbf{45 vanilla-NeRF runs} (9 pose-noise levels $\times$ 5 methods)
        to isolate when joint camera-pose optimization helps or hurts reconstruction.
  \item Showed a camera learning-rate warm-up recovers most quality lost to noisy poses
        (\textbf{14.1 $\rightarrow$ 22.7\,dB} PSNR at 5\% translation noise).
  \item Found that enabling pose optimization from perfect poses drops PSNR from 27.1 to 17.9\,dB, and that a
        Sim(3) alignment fit on training cameras alone recovers \textbf{+5.5\,dB}, revealing global coordinate
        drift rather than a worse reconstruction.
  \item Released v1.0.0 with reproducible configs, CI (GitHub Actions) and \texttt{uv}-managed environments.
\end{itemize}

% TODO: add 1--2 more projects (course projects, undergraduate thesis at HIT, etc.)
% \project{Project Title}{[\href{https://github.com/}{code}]}{2025}
% \begin{itemize}
%   \item What you built, with which method, and a quantified result.
% \end{itemize}

%---------------- Experience ----------------
% Uncomment when you have internships / RA positions.
% \section{Experience}
% \entry{Company / Lab}{Jun 2026 -- Present}{Research Intern, Team Name}{City}
% \begin{itemize}
%   \item Verb + what + how + measurable impact.
% \end{itemize}

%---------------- Honors ----------------
% \section{Honors \& Awards}
% \begin{itemize}
%   \item Scholarship / competition award \hfill 2024  % TODO
% \end{itemize}

%---------------- Skills ----------------
\section{Skills}
\begin{tabularx}{\linewidth}{@{}l@{\hspace{1.2em}}X@{}}
  \textbf{Deep Learning \& Vision} & PyTorch, OpenCV, NumPy, NeRF \\
  \textbf{Programming}                  & Python, C/C++, MATLAB \\
  \textbf{Systems}                      & Linux, ROS \\
  \textbf{Dev Tools}                    & Git, GitHub Actions (CI), uv, Conda, LaTeX \\
  \textbf{Languages}                    & Mandarin (native), English (TODO: CET-6 / IELTS score) \\
\end{tabularx}

\end{document}
